\section{Identifying anti-competitive behaviors}\label{coll}
Our empirical strategy follows the approach of \citet{kawai2023using} to detect subtle signs of collusion or favoritism in auction outcomes. The core idea is to focus on procurement auctions where the competition is extremely tight, meaning the winner's bid is nearly identical to the runner-up's bid. Under genuinely competitive conditions, two firms bidding almost the same price should be similar in their cost efficiency and should not systematically differ in their past bidding success. Thus, absent any collusion or external favoritism, one would expect no abrupt differences in firm characteristics when moving from just losing an auction to just winning it.

We implement this idea using a regression discontinuity design (RDD) centered around the margin of victory. Let $\Delta_{i,t}$ be the difference between firm $i$'s bid and the lowest competing bid in auction $t$. This margin $\Delta_{i,t}$ serves as our running variable, with a cutoff at zero. Firms with $\Delta_{i,t} < 0$ win the auction (their bid is lower than the best alternative), whereas those with $\Delta_{i,t} > 0$ lose. Intuitively, $\Delta_{i,t} = 0$ represents a threshold between winning and losing. We focus on auctions where $|\Delta_{i,t}|$ is very small, meaning the winner barely underbid the closest rival. In such cases, any systematic jump in firm outcomes at $\Delta_{i,t}=0$ can be interpreted as evidence of an outside influence on the auction result.

The two primary outcomes we examine are the firm's \emph{incumbency status} and its bid timing, identifying whether it is the \emph{last bidder} in the auction. We define an \textit{incumbent} as a firm that won the last auction it participated in prior to the current auction (incumbency is an indicator variable equal to 1 if the firm was the previous winner and 0 otherwise). \textit{Last bidder} identifies whether a firm was the last among the competitors to submit the bid. By examining these variables for firms just on either side of the winning threshold, we can detect irregular patterns: for instance, an unexpectedly high incumbency rate or unusually larger probability of being the last bidder for firms that barely win compared to those that barely lose.

Formally, we estimate the discontinuity at $\Delta=0$ for these outcomes. Let $x_{i,t}$ denote a generic outcome (such as the incumbency indicator or backlog) for firm $i$ before auction $t$. The discontinuity is captured by the difference in the expected value of $x_{i,t}$ immediately to the left and right of the cutoff:
\[
\beta \;=\; \lim_{\epsilon \to 0^+} \mathbb{E}[\,x_{i,t}\mid \Delta_{i,t}=\epsilon\,] \;-\; 
\lim_{\epsilon \to 0^-} \mathbb{E}[\,x_{i,t}\mid \Delta_{i,t}=\epsilon\,]~,
\]
where $\epsilon$ is a small positive increment in the bid margin. In the absence of any strategic behavior or external interference, we would expect $\beta \approx 0$ for characteristics like incumbency, since firms just winning or just losing should be comparable. Conversely, a significant $\beta \neq 0$ indicates a discontinuous jump in firm characteristics at the threshold, consistent with non-competitive forces affecting the outcome.

We hypothesize that a positive advantage conferred to narrowly winning firms would manifest as $\beta > 0$ in our setup (because the outcome for winners exceeds that for losers). Detecting a substantial gap in these measures at $\Delta=0$ will allow us to distinguish between a competitive market and one influenced by collusion or favoritism. 


For example, if there is political favoritism benefiting certain firms, those firms would appear as incumbents much more often when they win narrowly, whereas otherwise identical firms that narrowly lose would not show such a history. Similarly, a collusive bid-rotation scheme might lead to a more equal distribution of recent wins among participants, in which case we would \emph{not} expect a sharp drop in backlog for losing firms at the threshold. 
